% Lösungen Einheit 3 von 4 – Klammern auflösen (zusammenbau v0.9)
\einheitenkopf{Einheit 3 von 4 · Klammern auflösen}

\erg{27}{a) ein Plus \quad b) ein Minus \quad c) die Zahl $5$ als Faktor \quad d) $-6$ als Faktor}
\erg{28}{a) Klammer weglassen: $5a + (2b - 1) = 5a + 2b - 1$ \quad b) Klammer weglassen: $5a + (2b - 3) = 5a + 2b - 3$ \quad c) Klammer weglassen: $5a + (2b - 4) = 5a + 2b - 4$ \quad d) Klammer weglassen: $5a + (2b - 6) = 5a + 2b - 6$ \quad e) $5x + 15$ \quad f) Klammer zuerst: $30 - 20 = 10$; Glied für Glied: $30 - 12 - 8 = 10$; Ergebnis: beide Wege ergeben 10.}
\erg{29}{a) $5a - 2b - 4$ \quad b) $7a - 2b + 5 - c$ \quad c) $-2x - 10$ \quad d) $7x + 12$}
\erg{30}{$4x + 23$}
\erg{31}{a) Das Minus vor der Klammer dreht beide Vorzeichen, aus $-5$ wird $+5$. Richtig: Klammer auflösen: $14 - (a - 5) = 14 - a + 5$; zusammenfassen: $= 19 - a$. \quad b) Richtig. $9b - (4b - 2) = 9b - 4b + 2 = 5b + 2$; das Minus vor der Klammer dreht alle Vorzeichen in der Klammer, aus $-2$ wird $+2$. \quad c) Nicht stimmen können die Rechnungen 2 und 3. (2): Die Zahl vor der Klammer muss mit jedem Glied malgenommen werden, hinten muss ein Vielfaches von 3 stehen. (3): Das Minus vor der Klammer dreht alle Vorzeichen, aus $-6$ muss $+6$ werden.}
\erg{32}{a) Ja: $5 - (x - 2) = 5 - x + 2 = 7 - x$. \quad b) Nein; das Minus vor der Klammer dreht alle Vorzeichen, auch das der 3: $15 - (x - 3) = 15 - x + 3 = 18 - x$. \quad c) (1) wahr, denn Plus vor der Klammer ändert kein Vorzeichen. (2) falsch, z. B. $10 - (x + 1) = 10 - x - 1$: aus $+x$ wird $-x$. (3) wahr, denn so sagt es das Distributivgesetz.}
\erg{33}{a) $5 \cdot (x + 2) = 5x + 10$ \quad b) $2 \cdot (a + 6) = 2a + 12$ \quad c) Fehlende Einträge $+7$ und $4x$; $4 \cdot (x + 7) = 4x + 28$}
\erg{34}{a) Term mit Klammer: $9 \cdot (24 + x)$; Klammer auflösen: $9 \cdot (24 + x) = 216 + 9x$. Die Klasse zahlt $216 + 9x$ Euro. \quad b) Nein; Term mit Klammer: $2 \cdot (a + 12)$; Klammer auflösen: $2a + 24$; einsetzen: $2 \cdot 20 + 24 = 64$ cm, das sind 4 cm mehr als 60 cm. \quad c) Term mit Klammer: $2{,}50 \cdot (20 + x)$; Klammer auflösen: $50 + 2{,}50x$; einsetzen: $50 + 2{,}50 \cdot 4 = 60$ €. Die Muffins kosten 60 €.}
